\chapter{The event data model}
\label{ch:datamodel}

\chaptergoals{5 of 11 in Part~II.}
{\cref{ch:bigpicture}; the HNL input file from \cref{ch:setup}.}
{name the event classes and their units, describe the layout of each kind of file, and open one from Python.}

\section{Plain structs, written as RNTuples}

\repo{data-model} defines the classes every package exchanges. They are
deliberately simple: plain C++ structs with default values, no methods to
speak of, no pointers between objects, and ROOT dictionaries so that they can
be written to and read from ROOT files. Relations are expressed by
\emph{indices} (a mother index into the same particle vector, a Geant4 track
ID), never by pointers, which keeps the files readable from any language.

The files are written as \href{https://root.cern/doc/latest/classROOT_1_1RNTuple.html}{RNTuples},
ROOT's newer columnar I/O format, designed for efficient reading and writing, with
with a stable on-disk format since ROOT~6.34. An RNTuple has named
\emph{fields}; a field can be a number, a struct or a \texttt{std::vector} of
structs, and every member of a struct is itself a (sub)field.

Dimensional floating-point quantities use \texttt{double}; IDs and statuses
are integers. The storage units are:
\textbf{mm, GeV, ns}. The headers also provide typed views
(\file{SHiP/QuantityView.hpp}) that attach physical units at compile time;
\pkg{aegir} uses them internally, but you do not need them to read a file.

\section{The classes in this tutorial}

\begin{center}
\small
\begin{tabularx}{\linewidth}{P{2.6cm} P{2.9cm} L}
\toprule
Class & Produced by & Content \\
\midrule
\texttt{EventHeader}   & generator / converter & event \field{weight}, \field{original\_event\_id} \\
\texttt{MCParticle}    & generator / converter & generator record: \field{pdgCode}, \field{vertex}, \field{momentum},
                                                \field{energy} (total), \field{time}, \field{motherId},
                                                \field{mothers}, \field{status} \\
\texttt{SimParticle}   & \pkg{aegir} (Geant4)  & every tracked particle: \field{trackId}, \field{parentId},
                                                \field{pdgCode}, \field{vertex}, \field{endpoint},
                                                \field{momentum}, \field{energy} (kinetic),
                                                \field{time}, \field{creatorProcess} \\
\texttt{SimHit}        & \pkg{aegir} (Geant4)  & one positive-energy-deposit step in a sensitive volume (default scoring mode): \field{detectorId},
                                                \field{geometryNodeId}, \field{trackId}, \field{pdgCode},
                                                \field{position}, \field{momentum},
                                                \field{energyDeposit}, \field{time}, \field{pathLength} \\
\texttt{RecHit}        & \pkg{Shannon}         & a digitised hit; same members as \texttt{SimHit}
                                                minus \field{geometryNodeId} \\
detector hits          & \pkg{Shannon}         & \texttt{UBTHit}, \texttt{SBTHit}, \texttt{StrawTubesHit},
                                                \texttt{CaloHit}, \texttt{TimeDetHit}: one per detector,
                                                each wrapping a \texttt{RecHit} as \field{recHit} \\
\bottomrule
\end{tabularx}
\end{center}

\Cref{app:datamodel} lists every member with its unit, and the
\texttt{RecParticle} and \texttt{TrackFitResult} classes used by
reconstruction.

\subsection{Generator particles}

\texttt{MCParticle::status} follows the HepMC convention. Generator adapters
must map their native status codes to it; Pythia\,8 internal status codes
should not be copied unchanged. What matters here:

\begin{center}
\begin{tabular}{cl}
\toprule
\field{status} & meaning \\
\midrule
1 & final state: not decayed by the generator. \textbf{Only these are given to Geant4.} \\
2 & decayed by the generator \\
4 & incoming beam particle \\
\bottomrule
\end{tabular}
\end{center}

Mother links point into the same vector: \field{motherId} is the index of the
first mother, or $-1$; \field{mothers} lists all of them. For our HNL events
the vector is short, as \cref{ch:generator} shows: the HNL with status~2 and
its two daughters with status~1, pointing back to it.

\begin{pitfall}
\texttt{MCParticle::energy} is the \emph{total} energy;
\texttt{SimParticle::energy} is the \emph{kinetic} energy. Compute invariant
masses from \texttt{MCParticle}s, or recompute $E$ from $|\vec p|$ and the
mass.
\end{pitfall}

\subsection{Simulated particles and hits}

Every particle Geant4 tracks becomes a \texttt{SimParticle}. Its
\field{trackId} is the Geant4 track ID, unique \emph{within one event}; the
particles that came from the generator have \field{parentId}~$=0$, and every
secondary points to the track that created it. \field{creatorProcess} is the
Geant4 process sub-type that created it.

In the default scoring mode, a Geant4 step inside a sensitive volume becomes
a \texttt{SimHit} only if its deposited energy is positive. Crossing mode
records volume entrances and has a different meaning. Its
\field{trackId} says which particle made it, and its \field{detectorId} says
which detector it is in. The detector numbering is fixed by the data model:

\begin{center}
\begin{tabular}{cl}
\toprule
\field{detectorId} & detector \\
\midrule
0 & upstream background tagger (UBT) \\
1 & surround background tagger (SBT) \\
2 & straw tubes \\
3 & calorimeter \\
4 & timing detector \\
\bottomrule
\end{tabular}
\end{center}

\begin{pitfall}
\pkg{aegir} assigns \field{detectorId} as the \emph{position} of the matching
name in its \texttt{sensitive\_volumes} list, while \pkg{Shannon} interprets it
with the table above. An ID outside 0--4 raises ``No digitiser registered
for detector ID''; a wrong ordering within 0--4 silently calls the wrong digitiser. The list in the simulation workflow must
therefore have exactly five entries, in this order (\cref{sec:sensitive}).
\end{pitfall}

\section{What the files look like}
\label{sec:filelayout}

Three kinds of file appear in this tutorial. All of them hold their data in
RNTuples inside an ordinary ROOT file.

\paragraph{Generator input} (\file{hnl\_events.root}, from the converter).
One RNTuple, \texttt{events}, one entry per event:
\begin{center}\small
\begin{tabular}{ll}
\toprule
field & type \\
\midrule
\texttt{mc\_particles} & \texttt{std::vector<SHiP::MCParticle>} \\
\texttt{event\_header} & \texttt{SHiP::EventHeader} \\
\bottomrule
\end{tabular}
\end{center}

\paragraph{Simulation output} (\file{hnl\_sim.root}, from \pkg{aegir}).
One RNTuple, \texttt{events}, one entry per event:
\begin{center}\small
\begin{tabular}{ll}
\toprule
field & type \\
\midrule
\texttt{event\_number}  & \texttt{std::uint64\_t} \\
\texttt{event\_header}  & \texttt{SHiP::EventHeader} \\
\texttt{mc\_particles}  & \texttt{std::vector<SHiP::MCParticle>} \\
\texttt{sim\_hits}      & \texttt{std::vector<SHiP::SimHit>} \\
\texttt{sim\_particles} & \texttt{std::vector<SHiP::SimParticle>} \\
\bottomrule
\end{tabular}
\end{center}

\paragraph{Digitisation output} (\file{hnl\_digi.root}, from \pkg{Shannon}).
Seven RNTuples, one per collection, each with a single field \texttt{hit} and
\emph{one entry per object} (a hit, or a particle for \texttt{sim\_particles}):
\texttt{ubt\_hits}, \texttt{sbt\_hits}, \texttt{straw\_tubes\_hits},
\texttt{calorimeter\_hits}, \texttt{timing\_detector\_hits}, and copies of
the input as \texttt{sim\_hits} and \texttt{sim\_particles}.
\Cref{ch:digitisation} discusses what this layout means for an analysis.

\begin{exercise}{Look inside the HNL input file}
\label{ex:openfile}
In \file{work/}, start Python in the \pkg{aegir} environment and open the
file you were given:
\begin{py}
import ROOT
reader = ROOT.RNTupleReader.Open("events", "hnl_events.root")
reader.PrintInfo()     # fields, types, number of entries, compression
reader.Show(0)         # the first entry, as JSON
\end{py}
How many events does the file hold? How many particles are in the first
event, what are their PDG codes and statuses, and which ones will Geant4
track?
\end{exercise}

\begin{solution}
\texttt{Show(0)} lists one HNL (a converter-defined HNL PDG code and status~2)
and two daughters (PDG $(211,13)$ or $(-211,-13)$, status~1, \field{motherId}~$=0$).
Geant4 tracks only the two daughters. Opposite electric charges have the
same PDG-code sign here: $\pi^+$ is 211, while $\mu^-$ is 13. EventCalc's
text LLP column is a placeholder (12345678), so the converter must assign
and document the physical HNL code.
\end{solution}

\begin{note}
ROOT reads these files through the data-model dictionaries, which the pixi
environment makes available automatically; that is what lets you write
\texttt{p.pdgCode} on a particle in Python. Outside the environment,
\href{https://github.com/scikit-hep/uproot5}{uproot} can read the same
members column by column without any SHiP library.
\end{note}

\section{Keeping old files readable}

Backward readability is a project goal checked with frozen reference files
and a committed schema snapshot, but it has known exceptions. In 0.6,
reads of pre-0.5 fixtures are marked as expected failures for a ROOT issue.
The detector-wrapper rename \field{recHit} to \field{rec\_hit} has no working
conversion rule: reading an old wrapper intentionally raises a fatal error.
A successful read of a simulation fixture therefore does not demonstrate that
old Shannon output is readable. Keep the complete tutorial chain on 0.5
until an explicit migration and readback test succeeds. See
\file{data-model/tests/data/README.md} at the revision in \cref{app:review}.
